
\documentclass{article}
\usepackage{colortbl}
\usepackage{makecell}
\usepackage{multirow}
\usepackage{supertabular}

\begin{document}

\newcounter{utterance}

\centering \large Interaction Transcript for game `guesswhat', experiment `Level\_3', episode 0 with reka{-}flash{-}3.1{-}t0.0.
\vspace{24pt}

{ \footnotesize  \setcounter{utterance}{1}
\setlength{\tabcolsep}{0pt}
\begin{supertabular}{c@{$\;$}|p{.15\linewidth}@{}p{.15\linewidth}p{.15\linewidth}p{.15\linewidth}p{.15\linewidth}p{.15\linewidth}}
    \# & $\;$A & \multicolumn{4}{c}{Game Master} & $\;\:$B\\
    \hline

    \theutterance \stepcounter{utterance}  
    & & \multicolumn{4}{p{0.6\linewidth}}{
        \cellcolor[rgb]{0.9,0.9,0.9}{
            \makecell[{{p{\linewidth}}}]{
                \texttt{\tiny{[P1$\langle$GM]}}
                \texttt{Let's play a guessing game. Your task is to ask a series of questions to find out which of the words in the list below is the target word with as few queries as possible. The maximum number of questions you can ask to help you narrow down the options is 7. Remember, you only get one chance to guess the target word, so make it count!} \\
\\ 
\texttt{The rules of the game are: } \\
\\ 
\texttt{1. Only ask one question per turn. DO NOT ASK MORE THAN ONE QUESTION.} \\
\\ 
\texttt{2. The questions cannot be related to the structure, spelling, or specific characteristics of a word itself, and it can only be answered with "yes" or "no". } \\
\\ 
\texttt{3. Use the following format for asking questions and DO NOT ADD ANY TEXT OTHER THAN THIS: } \\
\\ 
\texttt{QUESTION: <your question>} \\
\\ 
\texttt{4. You have only one opportunity to guess the target word. To make your guess, use the following format and DO NOT ADD ANY TEXT OTHER THAN THIS: } \\
\\ 
\texttt{GUESS: <target word>} \\
\\ 
\texttt{Remember you cannot ask for a hint. Make your question brief. Let’s start!} \\
\\ 
\texttt{List of candidate words: {[}'Banjo', 'Guitar', 'Drum', 'Cymbal', 'Bugle', 'Trombone', 'Mandolin', 'Harp'{]}} \\
            }
        }
    }
    & \\ \\

    \theutterance \stepcounter{utterance}  
    & \multicolumn{4}{p{0.6\linewidth}}{
        \cellcolor[rgb]{0.9,0.9,0.9}{
            \makecell[{{p{\linewidth}}}]{
                \texttt{\tiny{[P1$\rangle$GM]}}
                \texttt{<reasoning>} \\
\texttt{Let's see. I need to figure out which word from the list is the target using up to 7 yes/no questions. The list is {[}'Banjo', 'Guitar', 'Drum', 'Cymbal', 'Bugle', 'Trombone', 'Mandolin', 'Harp'{]}.} \\
\\ 
\texttt{First, I should categorize these words. Let's see: Banjo, Guitar, Mandolin are string instruments. Drum, Cymbal, Trombone, Bugle are percussion or brass. Harp is a string instrument. Wait, Trombone is a brass instrument, right? So maybe split into string and percussion/brass. But Harp is a string, so maybe first question: Is the target a string instrument? Let's check the list. Banjo, Guitar, Mandolin, Harp are strings. The others are percussion or brass. So if the answer is yes, then the target is among those four. If no, then it's Drum, Cymbal, Bugle, Trombone.} \\
\\ 
\texttt{But wait, Trombone is brass, so if the answer is no, the target is Drum, Cymbal, Bugle, Trombone. Let's see. Let's ask that first. So first question: Is the target a string instrument? If yes, then narrow down to Banjo, Guitar, Mandolin, Harp. If no, then Drum, Cymbal, Bugle, Trombone.} \\
\\ 
\texttt{Alternatively, maybe another approach. Let's think about the number of letters. But the rules say questions can't be about structure or spelling. So that's out. So maybe category{-}based questions.} \\
\\ 
\texttt{Another angle: Are any of these instruments typically played with hands? Like drums, cymbals, maybe trombone? Wait, trombone is played with hands and feet. Harp is definitely hands. Banjo, guitar, mandolin are plucked with fingers or picks. So maybe: Is the target played with hands? If yes, then Harp, Banjo, Guitar, Mandolin. If no, then Drum, Cymbal, Bugle, Trombone. Wait, but drums are played with sticks or hands. Hmm, maybe that's not a clear split. Alternatively, maybe think about the primary method of producing sound. String instruments use strings, percussion use vibration, brass use air. So maybe first question: Does the target produce sound through vibration (like drums, c} \\
            }
        }
    }
    & & \\ \\

    \theutterance \stepcounter{utterance}  
    & & & \multicolumn{2}{p{0.3\linewidth}}{
        \cellcolor[rgb]{0.9,0.9,0.9}{
            \makecell[{{p{\linewidth}}}]{
                \texttt{\tiny{[GM$|$GM]}}
                \texttt{Invalid format. Guesser must use the form 'QUESTION: ' or 'GUESS: '.} \\
            }
        }
    }
    & & \\ \\

    \theutterance \stepcounter{utterance}  
    & & & \multicolumn{2}{p{0.3\linewidth}}{
        \cellcolor[rgb]{0.9,0.9,0.9}{
            \makecell[{{p{\linewidth}}}]{
                \texttt{\tiny{[GM$|$GM]}}
                \texttt{abort game} \\
            }
        }
    }
    & & \\ \\

\end{supertabular}
}

\end{document}
